% Options for packages loaded elsewhere
\PassOptionsToPackage{unicode}{hyperref}
\PassOptionsToPackage{hyphens}{url}
\documentclass[
]{article}
\usepackage{xcolor}
\usepackage[margin=1in]{geometry}
\usepackage{amsmath,amssymb}
\setcounter{secnumdepth}{-\maxdimen} % remove section numbering
\usepackage{iftex}
\ifPDFTeX
  \usepackage[T1]{fontenc}
  \usepackage[utf8]{inputenc}
  \usepackage{textcomp} % provide euro and other symbols
\else % if luatex or xetex
  \usepackage{unicode-math} % this also loads fontspec
  \defaultfontfeatures{Scale=MatchLowercase}
  \defaultfontfeatures[\rmfamily]{Ligatures=TeX,Scale=1}
\fi
\usepackage{lmodern}
\ifPDFTeX\else
  % xetex/luatex font selection
\fi
% Use upquote if available, for straight quotes in verbatim environments
\IfFileExists{upquote.sty}{\usepackage{upquote}}{}
\IfFileExists{microtype.sty}{% use microtype if available
  \usepackage[]{microtype}
  \UseMicrotypeSet[protrusion]{basicmath} % disable protrusion for tt fonts
}{}
\makeatletter
\@ifundefined{KOMAClassName}{% if non-KOMA class
  \IfFileExists{parskip.sty}{%
    \usepackage{parskip}
  }{% else
    \setlength{\parindent}{0pt}
    \setlength{\parskip}{6pt plus 2pt minus 1pt}}
}{% if KOMA class
  \KOMAoptions{parskip=half}}
\makeatother
\usepackage{color}
\usepackage{fancyvrb}
\newcommand{\VerbBar}{|}
\newcommand{\VERB}{\Verb[commandchars=\\\{\}]}
\DefineVerbatimEnvironment{Highlighting}{Verbatim}{commandchars=\\\{\}}
% Add ',fontsize=\small' for more characters per line
\usepackage{framed}
\definecolor{shadecolor}{RGB}{248,248,248}
\newenvironment{Shaded}{\begin{snugshade}}{\end{snugshade}}
\newcommand{\AlertTok}[1]{\textcolor[rgb]{0.94,0.16,0.16}{#1}}
\newcommand{\AnnotationTok}[1]{\textcolor[rgb]{0.56,0.35,0.01}{\textbf{\textit{#1}}}}
\newcommand{\AttributeTok}[1]{\textcolor[rgb]{0.13,0.29,0.53}{#1}}
\newcommand{\BaseNTok}[1]{\textcolor[rgb]{0.00,0.00,0.81}{#1}}
\newcommand{\BuiltInTok}[1]{#1}
\newcommand{\CharTok}[1]{\textcolor[rgb]{0.31,0.60,0.02}{#1}}
\newcommand{\CommentTok}[1]{\textcolor[rgb]{0.56,0.35,0.01}{\textit{#1}}}
\newcommand{\CommentVarTok}[1]{\textcolor[rgb]{0.56,0.35,0.01}{\textbf{\textit{#1}}}}
\newcommand{\ConstantTok}[1]{\textcolor[rgb]{0.56,0.35,0.01}{#1}}
\newcommand{\ControlFlowTok}[1]{\textcolor[rgb]{0.13,0.29,0.53}{\textbf{#1}}}
\newcommand{\DataTypeTok}[1]{\textcolor[rgb]{0.13,0.29,0.53}{#1}}
\newcommand{\DecValTok}[1]{\textcolor[rgb]{0.00,0.00,0.81}{#1}}
\newcommand{\DocumentationTok}[1]{\textcolor[rgb]{0.56,0.35,0.01}{\textbf{\textit{#1}}}}
\newcommand{\ErrorTok}[1]{\textcolor[rgb]{0.64,0.00,0.00}{\textbf{#1}}}
\newcommand{\ExtensionTok}[1]{#1}
\newcommand{\FloatTok}[1]{\textcolor[rgb]{0.00,0.00,0.81}{#1}}
\newcommand{\FunctionTok}[1]{\textcolor[rgb]{0.13,0.29,0.53}{\textbf{#1}}}
\newcommand{\ImportTok}[1]{#1}
\newcommand{\InformationTok}[1]{\textcolor[rgb]{0.56,0.35,0.01}{\textbf{\textit{#1}}}}
\newcommand{\KeywordTok}[1]{\textcolor[rgb]{0.13,0.29,0.53}{\textbf{#1}}}
\newcommand{\NormalTok}[1]{#1}
\newcommand{\OperatorTok}[1]{\textcolor[rgb]{0.81,0.36,0.00}{\textbf{#1}}}
\newcommand{\OtherTok}[1]{\textcolor[rgb]{0.56,0.35,0.01}{#1}}
\newcommand{\PreprocessorTok}[1]{\textcolor[rgb]{0.56,0.35,0.01}{\textit{#1}}}
\newcommand{\RegionMarkerTok}[1]{#1}
\newcommand{\SpecialCharTok}[1]{\textcolor[rgb]{0.81,0.36,0.00}{\textbf{#1}}}
\newcommand{\SpecialStringTok}[1]{\textcolor[rgb]{0.31,0.60,0.02}{#1}}
\newcommand{\StringTok}[1]{\textcolor[rgb]{0.31,0.60,0.02}{#1}}
\newcommand{\VariableTok}[1]{\textcolor[rgb]{0.00,0.00,0.00}{#1}}
\newcommand{\VerbatimStringTok}[1]{\textcolor[rgb]{0.31,0.60,0.02}{#1}}
\newcommand{\WarningTok}[1]{\textcolor[rgb]{0.56,0.35,0.01}{\textbf{\textit{#1}}}}
\usepackage{graphicx}
\makeatletter
\newsavebox\pandoc@box
\newcommand*\pandocbounded[1]{% scales image to fit in text height/width
  \sbox\pandoc@box{#1}%
  \Gscale@div\@tempa{\textheight}{\dimexpr\ht\pandoc@box+\dp\pandoc@box\relax}%
  \Gscale@div\@tempb{\linewidth}{\wd\pandoc@box}%
  \ifdim\@tempb\p@<\@tempa\p@\let\@tempa\@tempb\fi% select the smaller of both
  \ifdim\@tempa\p@<\p@\scalebox{\@tempa}{\usebox\pandoc@box}%
  \else\usebox{\pandoc@box}%
  \fi%
}
% Set default figure placement to htbp
\def\fps@figure{htbp}
\makeatother
\setlength{\emergencystretch}{3em} % prevent overfull lines
\providecommand{\tightlist}{%
  \setlength{\itemsep}{0pt}\setlength{\parskip}{0pt}}
\usepackage{bookmark}
\IfFileExists{xurl.sty}{\usepackage{xurl}}{} % add URL line breaks if available
\urlstyle{same}
\hypersetup{
  pdftitle={Case Study 2},
  pdfauthor={Maximilian Kraus, Lukas Vater, Hannah Weichhart},
  hidelinks,
  pdfcreator={LaTeX via pandoc}}

\title{Case Study 2}
\author{Maximilian Kraus, Lukas Vater, Hannah Weichhart}
\date{2026-04-02}

\begin{document}
\maketitle

\emph{The .Rmd} \textbf{and} \emph{.html (or .pdf) should be uploaded in
TUWEL by the deadline. Refrain from using explanatory comments in the R
code chunks but write them as text instead. Points will be deducted if
the submitted file is not in a decent form.}

\textbf{DISCLAIMER}: In case students did not contribute equally,
include a disclaimer stating what each student's contribution was.

The CIA World Factbook provides intelligence on various aspects of 266
world entities, including history, people, government, economy, energy,
geography, environment, communications, transportation, military,
terrorism, and transnational issues. This case study involves analyzing
world data from 2020, focusing on:

\begin{itemize}
\tightlist
\item
  \textbf{Education Expenditure (\% of GDP)}
\item
  \textbf{Youth Unemployment Rate (15-24 years)}
\item
  \textbf{Net Migration Rate} (difference between the number of people
  entering and leaving a country per 1,000 persons)
\end{itemize}

The data was sourced from the
\href{https://www.cia.gov/the-world-factbook/about/archives/}{CIA World
Factbook Archives}. You are required to use \texttt{dplyr} for data
manipulation, while any package can be used for importing data.

\section{Tasks:}\label{tasks}

\subsection{a. Data Import and
Cleaning}\label{a.-data-import-and-cleaning}

Load the following datasets from TUWEL and ensure that missing values
are handled correctly and column names are clear. Each dataset should
ultimately contain only two columns: \textbf{country} and the respective
variable. Note that some data sets also contain information on the year
when the value was last updated.

\begin{itemize}
\item
  \texttt{rawdata\_369.txt} which contains the (estimated) public
  expenditure on education as a percent of GDP. \emph{Pay attention! The
  delimiter is 2 or more white spaces (one space would not work as it
  would separate country names which contain a space); you have to skip
  the first two lines}.
\item
  \texttt{rawdata\_373.csv} which contains the (estimated) youth
  unemployment rate (15-24) per country
\item
  \texttt{rawdata\_347.txt} which contains (estimated) net migration
  rate per country.
\end{itemize}

\begin{Shaded}
\begin{Highlighting}[]
\FunctionTok{library}\NormalTok{(tidyverse)}
\end{Highlighting}
\end{Shaded}

\begin{Shaded}
\begin{Highlighting}[]
\NormalTok{education }\OtherTok{\textless{}{-}}\NormalTok{ readr}\SpecialCharTok{::}\FunctionTok{read\_delim}\NormalTok{(}\StringTok{\textquotesingle{}rawdata\_369.txt\textquotesingle{}}\NormalTok{, }\AttributeTok{delim =} \StringTok{\textquotesingle{}}\SpecialCharTok{\textbackslash{}n}\StringTok{\textquotesingle{}}\NormalTok{, }\AttributeTok{skip =} \DecValTok{1}\NormalTok{, }\AttributeTok{col\_names =}\NormalTok{ F)}\SpecialCharTok{\%\textgreater{}\%}
\FunctionTok{separate}\NormalTok{(X1, }\AttributeTok{into=} \FunctionTok{c}\NormalTok{(}\StringTok{\textquotesingle{}rank\textquotesingle{}}\NormalTok{, }\StringTok{\textquotesingle{}country\textquotesingle{}}\NormalTok{, }\StringTok{\textquotesingle{}education\_expenditures\textquotesingle{}}\NormalTok{, }\StringTok{\textquotesingle{}year\textquotesingle{}}\NormalTok{), }\AttributeTok{sep =} \StringTok{\textquotesingle{}}\SpecialCharTok{\textbackslash{}\textbackslash{}}\StringTok{s\{2,\}\textquotesingle{}}\NormalTok{, }\AttributeTok{convert =}\NormalTok{ T)}\SpecialCharTok{\%\textgreater{}\%}
\FunctionTok{select}\NormalTok{(country, education\_expenditures)}

\NormalTok{unemployment }\OtherTok{\textless{}{-}}\NormalTok{ readr}\SpecialCharTok{::}\FunctionTok{read\_csv}\NormalTok{(}\StringTok{\textquotesingle{}rawdata\_373.csv\textquotesingle{}}\NormalTok{)}

\NormalTok{migration }\OtherTok{\textless{}{-}}\NormalTok{ readr}\SpecialCharTok{::}\FunctionTok{read\_delim}\NormalTok{(}\StringTok{\textquotesingle{}rawdata\_347.txt\textquotesingle{}}\NormalTok{, }\AttributeTok{delim =} \StringTok{\textquotesingle{}}\SpecialCharTok{\textbackslash{}n}\StringTok{\textquotesingle{}}\NormalTok{, }\AttributeTok{skip =} \DecValTok{1}\NormalTok{, }\AttributeTok{col\_names =}\NormalTok{ F)}\SpecialCharTok{\%\textgreater{}\%}
\FunctionTok{separate}\NormalTok{(X1, }\AttributeTok{into =} \FunctionTok{c}\NormalTok{(}\StringTok{\textquotesingle{}rank\textquotesingle{}}\NormalTok{, }\StringTok{\textquotesingle{}country\textquotesingle{}}\NormalTok{, }\StringTok{\textquotesingle{}net\_migration\_rate\textquotesingle{}}\NormalTok{, }\StringTok{\textquotesingle{}year\textquotesingle{}}\NormalTok{), }\AttributeTok{sep =} \StringTok{\textquotesingle{}}\SpecialCharTok{\textbackslash{}\textbackslash{}}\StringTok{s\{2,\}\textquotesingle{}}\NormalTok{, }\AttributeTok{convert =}\NormalTok{ T)}\SpecialCharTok{\%\textgreater{}\%}
\FunctionTok{select}\NormalTok{(country, net\_migration\_rate)}
\end{Highlighting}
\end{Shaded}

\subsection{b. Merging Raw Data}\label{b.-merging-raw-data}

Merge the datasets using \texttt{dplyr} on a unique key and retain the
union of all observations.

\begin{itemize}
\tightlist
\item
  What key are you using for merging? The country column is used as key
  for merging as it is the only one present in all three datasets.
\item
  Return the dimensions of the merged dataset.
\end{itemize}

\begin{Shaded}
\begin{Highlighting}[]
\NormalTok{merged\_ds }\OtherTok{\textless{}{-}} \FunctionTok{full\_join}\NormalTok{(education, migration, }\AttributeTok{by =} \StringTok{\textquotesingle{}country\textquotesingle{}}\NormalTok{)}\SpecialCharTok{\%\textgreater{}\%}
\FunctionTok{full\_join}\NormalTok{(unemployment, }\AttributeTok{by =} \FunctionTok{c}\NormalTok{(}\StringTok{\textquotesingle{}country\textquotesingle{}} \OtherTok{=} \StringTok{\textquotesingle{}country\_name\textquotesingle{}}\NormalTok{))}

\FunctionTok{head}\NormalTok{(merged\_ds)}
\end{Highlighting}
\end{Shaded}

\begin{verbatim}
## # A tibble: 6 x 4
##   country            education_expenditures net_migration_rate youth_unempl_rate
##   <chr>                               <dbl>              <dbl>             <dbl>
## 1 Cuba                                 12.8               -3.7               6.1
## 2 Micronesia, Feder~                   12.5              -20.9              18.9
## 3 Solomon Islands                       9.9               -1.6               1.3
## 4 Montserrat                            8.8                0                NA  
## 5 Norway                                7.9                4                 9.7
## 6 Denmark                               7.8                2.8               9.4
\end{verbatim}

The dimension is: 227, 4, i.e.~the three observations for 227 countries.

\subsection{c.~Enriching Data with Income
Classification}\label{c.-enriching-data-with-income-classification}

Obtain country income classification (low, lower-middle, upper-middle,
high) from the
\href{https://datahelpdesk.worldbank.org/knowledgebase/articles/906519}{World
Bank} and merge it with the dataset.

\begin{itemize}
\tightlist
\item
  Identify common variables between datasets. Can they be used for
  merging? Why or why not? A reasonable choice would we to merge by
  country name. But since the data stems from different sources, it can
  not be ensured that the names are all spelled the same. Two examples
  for both the Democratic Republic of Congo and the Republic of Congo
  can be found down below. The merge should have resulted in two instead
  of four rows.
\end{itemize}

\begin{Shaded}
\begin{Highlighting}[]
\NormalTok{income\_class}\OtherTok{\textless{}{-}}\NormalTok{readr}\SpecialCharTok{::}\FunctionTok{read\_delim}\NormalTok{(}\StringTok{\textquotesingle{}Country\_Analytical\_History.csv\textquotesingle{}}\NormalTok{, }\AttributeTok{skip =} \DecValTok{11}\NormalTok{, }\AttributeTok{delim =} \StringTok{\textquotesingle{};\textquotesingle{}}\NormalTok{, }\AttributeTok{col\_names =} \FunctionTok{c}\NormalTok{(}\StringTok{\textquotesingle{}iso\textquotesingle{}}\NormalTok{, }\StringTok{\textquotesingle{}country\textquotesingle{}}\NormalTok{, }\StringTok{\textquotesingle{}1987\textquotesingle{}}\SpecialCharTok{:}\StringTok{\textquotesingle{}2024\textquotesingle{}}\NormalTok{))}\SpecialCharTok{\%\textgreater{}\%}
  \FunctionTok{select}\NormalTok{(iso, country, }\StringTok{\textquotesingle{}income\_class\textquotesingle{}} \OtherTok{=} \StringTok{\textquotesingle{}2020\textquotesingle{}}\NormalTok{)}

\NormalTok{full\_ds }\OtherTok{\textless{}{-}} \FunctionTok{full\_join}\NormalTok{(merged\_ds, income\_class, }\AttributeTok{by =} \StringTok{\textquotesingle{}country\textquotesingle{}}\NormalTok{)}

\FunctionTok{filter}\NormalTok{(full\_ds, }\FunctionTok{grepl}\NormalTok{(}\StringTok{\textquotesingle{}Congo\textquotesingle{}}\NormalTok{, country))}
\end{Highlighting}
\end{Shaded}

\begin{verbatim}
## # A tibble: 4 x 6
##   country      education_expenditures net_migration_rate youth_unempl_rate iso  
##   <chr>                         <dbl>              <dbl>             <dbl> <chr>
## 1 Congo, Repu~                    3.5               -0.9              NA   <NA> 
## 2 Congo, Demo~                    1.5               -0.9               8.7 <NA> 
## 3 Congo, Dem.~                   NA                 NA                NA   COD  
## 4 Congo, Rep.                    NA                 NA                NA   COG  
## # i 1 more variable: income_class <chr>
\end{verbatim}

\begin{itemize}
\tightlist
\item
  Since ISO codes are standardized, download and use the
  \href{https://www.cia.gov/the-world-factbook/references/country-data-codes/}{CIA
  country data codes} for merging. Make sure you are not losing any of
  the countries in your original data set when merging.
\end{itemize}

The names of Turkey and North Macedonia needed to be manually adjusted,
as they are inconsistent across the datasets.

\begin{Shaded}
\begin{Highlighting}[]
\NormalTok{country\_codes }\OtherTok{\textless{}{-}}\NormalTok{ readr}\SpecialCharTok{::}\FunctionTok{read\_delim}\NormalTok{(}\StringTok{\textquotesingle{}country\_data\_codes.csv\textquotesingle{}}\NormalTok{)}
\NormalTok{country\_codes }\OtherTok{\textless{}{-}} \FunctionTok{select}\NormalTok{(country\_codes, }\FunctionTok{c}\NormalTok{(Name, GENC))}
\FunctionTok{head}\NormalTok{(country\_codes)}
\end{Highlighting}
\end{Shaded}

\begin{verbatim}
## # A tibble: 6 x 2
##   Name           GENC 
##   <chr>          <chr>
## 1 Afghanistan    AFG  
## 2 Akrotiri       XQZ  
## 3 Albania        ALB  
## 4 Algeria        DZA  
## 5 American Samoa ASM  
## 6 Andorra        AND
\end{verbatim}

\begin{Shaded}
\begin{Highlighting}[]
\NormalTok{coded\_ds }\OtherTok{\textless{}{-}}\NormalTok{ merged\_ds}\SpecialCharTok{\%\textgreater{}\%}
  \FunctionTok{mutate}\NormalTok{(}\AttributeTok{country =} \FunctionTok{ifelse}\NormalTok{(country }\SpecialCharTok{==} \StringTok{"Macedonia"}\NormalTok{, }\StringTok{"North Macedonia"}\NormalTok{, country))}\SpecialCharTok{\%\textgreater{}\%}
  \FunctionTok{mutate}\NormalTok{(}\AttributeTok{country =} \FunctionTok{ifelse}\NormalTok{(country }\SpecialCharTok{==} \StringTok{"Turkey"}\NormalTok{, }\StringTok{"Turkey (Turkiye)"}\NormalTok{, country))}\SpecialCharTok{\%\textgreater{}\%}
  \FunctionTok{left\_join}\NormalTok{(country\_codes, }\AttributeTok{by =} \FunctionTok{c}\NormalTok{(}\StringTok{\textquotesingle{}country\textquotesingle{}} \OtherTok{=} \StringTok{\textquotesingle{}Name\textquotesingle{}}\NormalTok{))}\SpecialCharTok{\%\textgreater{}\%}
  \FunctionTok{left\_join}\NormalTok{(}\FunctionTok{select}\NormalTok{(income\_class, }\FunctionTok{c}\NormalTok{(iso, income\_class)), }\AttributeTok{by =} \FunctionTok{c}\NormalTok{(}\StringTok{\textquotesingle{}GENC\textquotesingle{}} \OtherTok{=} \StringTok{\textquotesingle{}iso\textquotesingle{}}\NormalTok{))}\SpecialCharTok{\%\textgreater{}\%}
  \FunctionTok{arrange}\NormalTok{(country)}

\FunctionTok{head}\NormalTok{(coded\_ds)}
\end{Highlighting}
\end{Shaded}

\begin{verbatim}
## # A tibble: 6 x 6
##   country      education_expenditures net_migration_rate youth_unempl_rate GENC 
##   <chr>                         <dbl>              <dbl>             <dbl> <chr>
## 1 Afghanistan                     4.1               -0.1              17.6 AFG  
## 2 Albania                         3.6               -3.3              31.9 ALB  
## 3 Algeria                        NA                 -0.9              39.3 DZA  
## 4 American Sa~                   NA                -26.1              NA   ASM  
## 5 Andorra                         3.2                0                NA   AND  
## 6 Angola                          3.4               -0.2              39.4 AGO  
## # i 1 more variable: income_class <chr>
\end{verbatim}

The dimension of the dataset is now: 227, 6, including the additional
columns ``iso'' and ``income\_class''.

The previous example of the Democratic Repulic and Republic of the Congo
now works just fine!

\begin{Shaded}
\begin{Highlighting}[]
\FunctionTok{filter}\NormalTok{(coded\_ds, }\FunctionTok{grepl}\NormalTok{(}\StringTok{\textquotesingle{}Congo\textquotesingle{}}\NormalTok{, country))}
\end{Highlighting}
\end{Shaded}

\begin{verbatim}
## # A tibble: 2 x 6
##   country      education_expenditures net_migration_rate youth_unempl_rate GENC 
##   <chr>                         <dbl>              <dbl>             <dbl> <chr>
## 1 Congo, Demo~                    1.5               -0.9               8.7 COD  
## 2 Congo, Repu~                    3.5               -0.9              NA   COG  
## # i 1 more variable: income_class <chr>
\end{verbatim}

\subsection{d.~Adding Geographical
Information}\label{d.-adding-geographical-information}

Introduce continent and subcontinent (or region) data for each country.

\begin{itemize}
\tightlist
\item
  Find and download an appropriate online resource.
\item
  Merge this information into the dataset, naming the final dataset
  \texttt{df\_vars}. Make sure you are not losing any of the countries
  in your original data set when merging.
\end{itemize}

Using
\href{https://ourworldindata.org/grapher/world-regions-according-to-un?tab=table}{World
region according to UN}

\begin{Shaded}
\begin{Highlighting}[]
\NormalTok{country\_continents }\OtherTok{\textless{}{-}}\NormalTok{ readr}\SpecialCharTok{::}\FunctionTok{read\_delim}\NormalTok{(}\StringTok{\textquotesingle{}world\_regions\_according\_to\_un.csv\textquotesingle{}}\NormalTok{)}
\NormalTok{country\_continents }\OtherTok{\textless{}{-}} \FunctionTok{select}\NormalTok{(country\_continents, }\FunctionTok{c}\NormalTok{(}\StringTok{"Code"}\NormalTok{, }\StringTok{"World regions according to UN"}\NormalTok{))}
\FunctionTok{head}\NormalTok{(country\_continents)}
\end{Highlighting}
\end{Shaded}

\begin{verbatim}
## # A tibble: 6 x 2
##   Code  `World regions according to UN`
##   <chr> <chr>                          
## 1 AFG   Asia (UN)                      
## 2 ALA   Europe (UN)                    
## 3 ALB   Europe (UN)                    
## 4 DZA   Africa (UN)                    
## 5 ASM   Oceania (UN)                   
## 6 AND   Europe (UN)
\end{verbatim}

\begin{Shaded}
\begin{Highlighting}[]
\NormalTok{df\_vars }\OtherTok{\textless{}{-}}\NormalTok{ coded\_ds}\SpecialCharTok{\%\textgreater{}\%}
  \FunctionTok{left\_join}\NormalTok{(country\_continents, }\AttributeTok{by =} \FunctionTok{c}\NormalTok{(}\StringTok{"GENC"} \OtherTok{=} \StringTok{"Code"}\NormalTok{))}\SpecialCharTok{\%\textgreater{}\%}
  \FunctionTok{arrange}\NormalTok{(country)}

\FunctionTok{head}\NormalTok{(df\_vars)}
\end{Highlighting}
\end{Shaded}

\begin{verbatim}
## # A tibble: 6 x 7
##   country      education_expenditures net_migration_rate youth_unempl_rate GENC 
##   <chr>                         <dbl>              <dbl>             <dbl> <chr>
## 1 Afghanistan                     4.1               -0.1              17.6 AFG  
## 2 Albania                         3.6               -3.3              31.9 ALB  
## 3 Algeria                        NA                 -0.9              39.3 DZA  
## 4 American Sa~                   NA                -26.1              NA   ASM  
## 5 Andorra                         3.2                0                NA   AND  
## 6 Angola                          3.4               -0.2              39.4 AGO  
## # i 2 more variables: income_class <chr>, `World regions according to UN` <chr>
\end{verbatim}

As the dimension: 227, 7 of the final dataset shows, the number of
countries represented is still the same.

\subsection{e. Data Tidiness and Summary
Statistics}\label{e.-data-tidiness-and-summary-statistics}

\begin{itemize}
\tightlist
\item
  Evaluate the tidiness of \texttt{df\_vars} (observational units,
  variables, fixed vs.~measured variables). Make adjustments to tidy the
  data, if necessary.
\end{itemize}

Each row is an observation/country, every row holds one variable. Rename
``World regions according to UN'' to just Continent and removing (UN) to
make the table more readable. GENC and continent are fixed variables and
are therefore placed at the beginning of the dataset.

\begin{Shaded}
\begin{Highlighting}[]
\NormalTok{df\_vars }\OtherTok{\textless{}{-}}\NormalTok{ df\_vars }\SpecialCharTok{\%\textgreater{}\%}
  \FunctionTok{rename}\NormalTok{(}\AttributeTok{continent =} \StringTok{\textasciigrave{}}\AttributeTok{World regions according to UN}\StringTok{\textasciigrave{}}\NormalTok{) }\SpecialCharTok{\%\textgreater{}\%}
  \FunctionTok{mutate}\NormalTok{(}\AttributeTok{continent =} \FunctionTok{gsub}\NormalTok{(}\StringTok{"}\SpecialCharTok{\textbackslash{}\textbackslash{}}\StringTok{s*}\SpecialCharTok{\textbackslash{}\textbackslash{}}\StringTok{(UN}\SpecialCharTok{\textbackslash{}\textbackslash{}}\StringTok{)"}\NormalTok{, }\StringTok{""}\NormalTok{, continent)) }\SpecialCharTok{\%\textgreater{}\%}
  \FunctionTok{select}\NormalTok{(country, GENC, continent, }\FunctionTok{everything}\NormalTok{())}
\FunctionTok{head}\NormalTok{(df\_vars)}
\end{Highlighting}
\end{Shaded}

\begin{verbatim}
## # A tibble: 6 x 7
##   country        GENC  continent education_expenditures net_migration_rate
##   <chr>          <chr> <chr>                      <dbl>              <dbl>
## 1 Afghanistan    AFG   Asia                         4.1               -0.1
## 2 Albania        ALB   Europe                       3.6               -3.3
## 3 Algeria        DZA   Africa                      NA                 -0.9
## 4 American Samoa ASM   Oceania                     NA                -26.1
## 5 Andorra        AND   Europe                       3.2                0  
## 6 Angola         AGO   Africa                       3.4               -0.2
## # i 2 more variables: youth_unempl_rate <dbl>, income_class <chr>
\end{verbatim}

\begin{itemize}
\tightlist
\item
  Create a frequency table for the income status variable and briefly
  interpret the results.
\end{itemize}

\begin{Shaded}
\begin{Highlighting}[]
\NormalTok{income\_status }\OtherTok{\textless{}{-}}\NormalTok{ df\_vars }\SpecialCharTok{\%\textgreater{}\%}
  \FunctionTok{count}\NormalTok{(income\_class) }\SpecialCharTok{\%\textgreater{}\%}
  \FunctionTok{mutate}\NormalTok{(}\AttributeTok{percentage =}\NormalTok{ n }\SpecialCharTok{/} \FunctionTok{sum}\NormalTok{(n) }\SpecialCharTok{*} \DecValTok{100}\NormalTok{) }\SpecialCharTok{\%\textgreater{}\%}
  \FunctionTok{rename}\NormalTok{(}\AttributeTok{absolute =}\NormalTok{ n)}
\NormalTok{income\_status}
\end{Highlighting}
\end{Shaded}

\begin{verbatim}
## # A tibble: 5 x 3
##   income_class absolute percentage
##   <chr>           <int>      <dbl>
## 1 H                  79      34.8 
## 2 L                  27      11.9 
## 3 LM                 54      23.8 
## 4 UM                 54      23.8 
## 5 <NA>               13       5.73
\end{verbatim}

The frequency table shows that the largest group of countries falls into
the high-income category (H), accounting for about 35\% of all
observations. Lower-middle income (LM) and upper-middle income (UM)
countries each represent approximately 24\%, indicating that a large
share of countries is concentrated in middle-income groups. Low-income
countries (L) make up about 12\% of the sample, while a small share
(around 6\%) has missing income classification (NA).

\begin{itemize}
\tightlist
\item
  Analyze the distribution of income status across continents by
  computing absolute and relative frequencies. Comment on the findings.
\end{itemize}

\begin{Shaded}
\begin{Highlighting}[]
\NormalTok{income\_continents }\OtherTok{\textless{}{-}}\NormalTok{ df\_vars }\SpecialCharTok{\%\textgreater{}\%}
  \FunctionTok{group\_by}\NormalTok{(continent) }\SpecialCharTok{\%\textgreater{}\%}
  \FunctionTok{count}\NormalTok{(income\_class) }\SpecialCharTok{\%\textgreater{}\%}
  \FunctionTok{mutate}\NormalTok{(}\AttributeTok{percentage =}\NormalTok{ n }\SpecialCharTok{/} \FunctionTok{sum}\NormalTok{(n) }\SpecialCharTok{*} \DecValTok{100}\NormalTok{) }\SpecialCharTok{\%\textgreater{}\%}
  \FunctionTok{rename}\NormalTok{(}\AttributeTok{absolute =}\NormalTok{ n)}

\NormalTok{income\_continents}\SpecialCharTok{\%\textgreater{}\%}
  \FunctionTok{pivot\_wider}\NormalTok{(}\SpecialCharTok{{-}}\NormalTok{percentage, }\AttributeTok{names\_from =}\NormalTok{ continent, }\AttributeTok{values\_from =}\NormalTok{ absolute)}
\end{Highlighting}
\end{Shaded}

\begin{verbatim}
## # A tibble: 5 x 8
##   income_class Africa  Asia Europe Latin America and the Ca~1 `Northern America`
##   <chr>         <int> <int>  <int>                      <int>              <int>
## 1 H                 1    14     35                         16                  4
## 2 L                23     4     NA                         NA                 NA
## 3 LM               23    18      1                          6                 NA
## 4 UM                7    13     10                         19                 NA
## 5 <NA>              1    NA      2                          4                  1
## # i abbreviated name: 1: `Latin America and the Caribbean`
## # i 2 more variables: Oceania <int>, `NA` <int>
\end{verbatim}

\begin{Shaded}
\begin{Highlighting}[]
\NormalTok{income\_continents}\SpecialCharTok{\%\textgreater{}\%}
  \FunctionTok{pivot\_wider}\NormalTok{(}\SpecialCharTok{{-}}\NormalTok{absolute, }\AttributeTok{names\_from =}\NormalTok{ continent, }\AttributeTok{values\_from =}\NormalTok{ percentage)}
\end{Highlighting}
\end{Shaded}

\begin{verbatim}
## # A tibble: 5 x 8
##   income_class Africa  Asia Europe Latin America and the Ca~1 `Northern America`
##   <chr>         <dbl> <dbl>  <dbl>                      <dbl>              <dbl>
## 1 H              1.82 28.6   72.9                       35.6                  80
## 2 L             41.8   8.16  NA                         NA                    NA
## 3 LM            41.8  36.7    2.08                      13.3                  NA
## 4 UM            12.7  26.5   20.8                       42.2                  NA
## 5 <NA>           1.82 NA      4.17                       8.89                 20
## # i abbreviated name: 1: `Latin America and the Caribbean`
## # i 2 more variables: Oceania <dbl>, `NA` <dbl>
\end{verbatim}

All low income countries are in Africa and Asia. Northern America has
only high income countries (and one NA). Europe consists overwhelmingly
of high and upper middle income countries, similar to Latin America and
the Caribbean, although in a less pronounced way. Oceania has a somewhat
more diverse structure, with more lower middle income countries than
those two continents.

\begin{itemize}
\tightlist
\item
  Using the distribution of income status across continents, identify
  which countries are the only ones in their income group across the
  continent. Discuss briefly.
\end{itemize}

\begin{Shaded}
\begin{Highlighting}[]
\NormalTok{unique\_groups }\OtherTok{\textless{}{-}}\NormalTok{ income\_continents }\SpecialCharTok{\%\textgreater{}\%}
  \FunctionTok{filter}\NormalTok{(absolute }\SpecialCharTok{==} \DecValTok{1}\NormalTok{)}

\NormalTok{unique\_countries }\OtherTok{\textless{}{-}}\NormalTok{ df\_vars }\SpecialCharTok{\%\textgreater{}\%}
  \FunctionTok{semi\_join}\NormalTok{(unique\_groups, }\AttributeTok{by =} \FunctionTok{c}\NormalTok{(}\StringTok{"continent"}\NormalTok{, }\StringTok{"income\_class"}\NormalTok{)) }\SpecialCharTok{\%\textgreater{}\%}
  \FunctionTok{select}\NormalTok{(country, continent, income\_class) }\SpecialCharTok{\%\textgreater{}\%}
  \FunctionTok{na.omit}\NormalTok{()}

\NormalTok{unique\_countries}
\end{Highlighting}
\end{Shaded}

\begin{verbatim}
## # A tibble: 2 x 3
##   country    continent income_class
##   <chr>      <chr>     <chr>       
## 1 Seychelles Africa    H           
## 2 Ukraine    Europe    LM
\end{verbatim}

Seychelles is the only high-income country in Africa, while Ukraine is
the only lower-middle-income country in Europe. Interestingly, Northern
America or Oceania do not have any single outlying countries, despite
the higher likelihood due to the small number of countries.

\subsection{f.~Further Summary Statistics and
Insights}\label{f.-further-summary-statistics-and-insights}

\begin{itemize}
\tightlist
\item
  Create a table of average (mean and median) values for expenditure,
  youth unemployment rate and net migration rate separated into income
  status. Make sure that in the output, the ordering of the income
  classes is proper (i.e., L, LM, UM, H or the other way around).
  Briefly comment the results and any differences between the mean and
  median.
\end{itemize}

\begin{Shaded}
\begin{Highlighting}[]
\NormalTok{summary\_table }\OtherTok{\textless{}{-}}\NormalTok{ df\_vars }\SpecialCharTok{\%\textgreater{}\%}
    \FunctionTok{mutate}\NormalTok{(}\AttributeTok{income\_class =} \FunctionTok{factor}\NormalTok{(income\_class, }\AttributeTok{levels =} \FunctionTok{c}\NormalTok{(}\StringTok{"H"}\NormalTok{, }\StringTok{"UM"}\NormalTok{, }\StringTok{"LM"}\NormalTok{, }\StringTok{"L"}\NormalTok{))) }\SpecialCharTok{\%\textgreater{}\%}
  \FunctionTok{group\_by}\NormalTok{(income\_class) }\SpecialCharTok{\%\textgreater{}\%}
  \FunctionTok{summarise}\NormalTok{(}
    \AttributeTok{mean\_edu\_expenditure =} \FunctionTok{mean}\NormalTok{(education\_expenditures, }\AttributeTok{na.rm =} \ConstantTok{TRUE}\NormalTok{),}
    \AttributeTok{median\_edu\_expenditure =} \FunctionTok{median}\NormalTok{(education\_expenditures, }\AttributeTok{na.rm =} \ConstantTok{TRUE}\NormalTok{),}
    
    \AttributeTok{mean\_unemployment =} \FunctionTok{mean}\NormalTok{(youth\_unempl\_rate, }\AttributeTok{na.rm =} \ConstantTok{TRUE}\NormalTok{),}
    \AttributeTok{median\_unemployment =} \FunctionTok{median}\NormalTok{(youth\_unempl\_rate, }\AttributeTok{na.rm =} \ConstantTok{TRUE}\NormalTok{),}
    
    \AttributeTok{mean\_migration =} \FunctionTok{mean}\NormalTok{(net\_migration\_rate, }\AttributeTok{na.rm =} \ConstantTok{TRUE}\NormalTok{),}
    \AttributeTok{median\_migration =} \FunctionTok{median}\NormalTok{(net\_migration\_rate, }\AttributeTok{na.rm =} \ConstantTok{TRUE}\NormalTok{)}
\NormalTok{  )}
\NormalTok{summary\_table}
\end{Highlighting}
\end{Shaded}

\begin{verbatim}
## # A tibble: 5 x 7
##   income_class mean_edu_expenditure median_edu_expenditure mean_unemployment
##   <fct>                       <dbl>                  <dbl>             <dbl>
## 1 H                            4.6                     4.7              16.5
## 2 UM                           4.36                    4.1              21.9
## 3 LM                           4.62                    4.2              16.5
## 4 L                            3.52                    3.1              14.7
## 5 <NA>                         6.02                    5.3              38.6
## # i 3 more variables: median_unemployment <dbl>, mean_migration <dbl>,
## #   median_migration <dbl>
\end{verbatim}

We can see that low middle income countries have higher relative
education expanditure and lower youth unemployment. Looking at the net
migration rate there is a big difference between mean and median,
especially for for upper middle income countries.

\begin{itemize}
\tightlist
\item
  Look at the standard deviation and the interquartile range of the
  variables per income status instead of the location statistics above.
  Do you gain additional insights? Briefly comment the results.
\end{itemize}

\begin{Shaded}
\begin{Highlighting}[]
\NormalTok{spread\_table }\OtherTok{\textless{}{-}}\NormalTok{ df\_vars }\SpecialCharTok{\%\textgreater{}\%}
  \FunctionTok{mutate}\NormalTok{(}\AttributeTok{income\_class =} \FunctionTok{factor}\NormalTok{(income\_class, }\AttributeTok{levels =} \FunctionTok{c}\NormalTok{(}\StringTok{"H"}\NormalTok{, }\StringTok{"UM"}\NormalTok{, }\StringTok{"LM"}\NormalTok{, }\StringTok{"L"}\NormalTok{))) }\SpecialCharTok{\%\textgreater{}\%}
  \FunctionTok{group\_by}\NormalTok{(income\_class) }\SpecialCharTok{\%\textgreater{}\%}
  \FunctionTok{summarise}\NormalTok{(}
    \AttributeTok{sd\_edu\_expenditure =} \FunctionTok{sd}\NormalTok{(education\_expenditures, }\AttributeTok{na.rm =} \ConstantTok{TRUE}\NormalTok{),}
    \AttributeTok{iqr\_edu\_expenditure =} \FunctionTok{IQR}\NormalTok{(education\_expenditures, }\AttributeTok{na.rm =} \ConstantTok{TRUE}\NormalTok{),}
    
    \AttributeTok{sd\_unemployment =} \FunctionTok{sd}\NormalTok{(youth\_unempl\_rate, }\AttributeTok{na.rm =} \ConstantTok{TRUE}\NormalTok{),}
    \AttributeTok{iqr\_unemployment =} \FunctionTok{IQR}\NormalTok{(youth\_unempl\_rate, }\AttributeTok{na.rm =} \ConstantTok{TRUE}\NormalTok{),}
    
    \AttributeTok{sd\_migration =} \FunctionTok{sd}\NormalTok{(net\_migration\_rate, }\AttributeTok{na.rm =} \ConstantTok{TRUE}\NormalTok{),}
    \AttributeTok{iqr\_migration =} \FunctionTok{IQR}\NormalTok{(net\_migration\_rate, }\AttributeTok{na.rm =} \ConstantTok{TRUE}\NormalTok{)}
\NormalTok{  )}
\NormalTok{spread\_table}
\end{Highlighting}
\end{Shaded}

\begin{verbatim}
## # A tibble: 5 x 7
##   income_class sd_edu_expenditure iqr_edu_expenditure sd_unemployment
##   <fct>                     <dbl>               <dbl>           <dbl>
## 1 H                          1.50                1.8             10.7
## 2 UM                         1.80                1.85            12.9
## 3 LM                         2.14                2.28            11.0
## 4 L                          1.69                2.4             12.7
## 5 <NA>                       1.87                1.03            17.2
## # i 3 more variables: iqr_unemployment <dbl>, sd_migration <dbl>,
## #   iqr_migration <dbl>
\end{verbatim}

There is a high standard deviation in migration for upper middle income
countries. Apart from migration, the spread appears to be relatively
similar across different income classes.

\begin{itemize}
\tightlist
\item
  Extend the analysis of the statistics median and IQR to \textbf{each
  income status and continent combination}. Play around with displaying
  the resulting table. Use \texttt{pivot\_longer()} and/or
  \texttt{pivot\_wider()} to generate different outputs. Discuss the
  results as well as the readability of the different tables.
\end{itemize}

\begin{Shaded}
\begin{Highlighting}[]
\NormalTok{table\_test }\OtherTok{\textless{}{-}}\NormalTok{ df\_vars }\SpecialCharTok{\%\textgreater{}\%}
  \FunctionTok{mutate}\NormalTok{(}\AttributeTok{income\_class =} \FunctionTok{factor}\NormalTok{(income\_class, }\AttributeTok{levels =} \FunctionTok{c}\NormalTok{(}\StringTok{"H"}\NormalTok{, }\StringTok{"UM"}\NormalTok{, }\StringTok{"LM"}\NormalTok{, }\StringTok{"L"}\NormalTok{))) }\SpecialCharTok{\%\textgreater{}\%}
  \FunctionTok{group\_by}\NormalTok{(continent, income\_class) }\SpecialCharTok{\%\textgreater{}\%}
  \FunctionTok{summarise}\NormalTok{(}
    \AttributeTok{median\_expenditure =} \FunctionTok{median}\NormalTok{(education\_expenditures, }\AttributeTok{na.rm =} \ConstantTok{TRUE}\NormalTok{),}
    \AttributeTok{iqr\_expenditure =} \FunctionTok{IQR}\NormalTok{(education\_expenditures, }\AttributeTok{na.rm =} \ConstantTok{TRUE}\NormalTok{),}
    
    \AttributeTok{median\_unemployment =} \FunctionTok{median}\NormalTok{(youth\_unempl\_rate, }\AttributeTok{na.rm =} \ConstantTok{TRUE}\NormalTok{),}
    \AttributeTok{iqr\_unemployment =} \FunctionTok{IQR}\NormalTok{(youth\_unempl\_rate, }\AttributeTok{na.rm =} \ConstantTok{TRUE}\NormalTok{),}
    
    \AttributeTok{median\_migration =} \FunctionTok{median}\NormalTok{(net\_migration\_rate, }\AttributeTok{na.rm =} \ConstantTok{TRUE}\NormalTok{),}
    \AttributeTok{iqr\_migration =} \FunctionTok{IQR}\NormalTok{(net\_migration\_rate, }\AttributeTok{na.rm =} \ConstantTok{TRUE}\NormalTok{),}
    \AttributeTok{.groups =} \StringTok{"drop"}
\NormalTok{  )}
\NormalTok{table\_test}
\end{Highlighting}
\end{Shaded}

\begin{verbatim}
## # A tibble: 25 x 8
##    continent income_class median_expenditure iqr_expenditure median_unemployment
##    <chr>     <fct>                     <dbl>           <dbl>               <dbl>
##  1 Africa    H                          4.4             0                   11.6
##  2 Africa    UM                         3.9             2.15                37  
##  3 Africa    LM                         4               1.9                 19.5
##  4 Africa    L                          2.95            2.55                 8.7
##  5 Africa    <NA>                      NA              NA                   NA  
##  6 Asia      H                          3.8             2.3                  8.9
##  7 Asia      UM                         3.1             1.45                18.0
##  8 Asia      LM                         4               2.3                 13.7
##  9 Asia      L                          4.1             0                   24.5
## 10 Europe    H                          4.8             1.5                 11.8
## # i 15 more rows
## # i 3 more variables: iqr_unemployment <dbl>, median_migration <dbl>,
## #   iqr_migration <dbl>
\end{verbatim}

\begin{Shaded}
\begin{Highlighting}[]
\NormalTok{table\_test\_2 }\OtherTok{\textless{}{-}}\NormalTok{ table\_test }\SpecialCharTok{\%\textgreater{}\%}
  \FunctionTok{pivot\_wider}\NormalTok{(}
    \AttributeTok{names\_from =}\NormalTok{ income\_class,}
    \AttributeTok{values\_from =} \FunctionTok{c}\NormalTok{(median\_expenditure, iqr\_expenditure,}
\NormalTok{                    median\_unemployment, iqr\_unemployment,}
\NormalTok{                    median\_migration, iqr\_migration)}
\NormalTok{  )}
\NormalTok{table\_test\_2}
\end{Highlighting}
\end{Shaded}

\begin{verbatim}
## # A tibble: 7 x 31
##   continent     median_expenditure_H median_expenditure_UM median_expenditure_LM
##   <chr>                        <dbl>                 <dbl>                 <dbl>
## 1 Africa                         4.4                   3.9                  4   
## 2 Asia                           3.8                   3.1                  4   
## 3 Europe                         4.8                   4.1                  5.4 
## 4 Latin Americ~                  4.9                   5                    5.25
## 5 Northern Ame~                  5                    NA                   NA   
## 6 Oceania                        5.7                   3.9                  4.5 
## 7 <NA>                          NA                    NA                   NA   
## # i 27 more variables: median_expenditure_L <dbl>, median_expenditure_NA <dbl>,
## #   iqr_expenditure_H <dbl>, iqr_expenditure_UM <dbl>,
## #   iqr_expenditure_LM <dbl>, iqr_expenditure_L <dbl>,
## #   iqr_expenditure_NA <dbl>, median_unemployment_H <dbl>,
## #   median_unemployment_UM <dbl>, median_unemployment_LM <dbl>,
## #   median_unemployment_L <dbl>, median_unemployment_NA <dbl>,
## #   iqr_unemployment_H <dbl>, iqr_unemployment_UM <dbl>, ...
\end{verbatim}

\begin{Shaded}
\begin{Highlighting}[]
\NormalTok{table\_test\_3 }\OtherTok{\textless{}{-}}\NormalTok{ table\_test }\SpecialCharTok{\%\textgreater{}\%}
  \FunctionTok{pivot\_wider}\NormalTok{(}
    \AttributeTok{names\_from =}\NormalTok{ continent,}
    \AttributeTok{values\_from =} \SpecialCharTok{{-}}\FunctionTok{c}\NormalTok{(continent, income\_class)}
\NormalTok{  )}
\NormalTok{table\_test\_3}
\end{Highlighting}
\end{Shaded}

\begin{verbatim}
## # A tibble: 5 x 43
##   income_class median_expenditure_Africa median_expenditure_Asia
##   <fct>                            <dbl>                   <dbl>
## 1 H                                 4.4                      3.8
## 2 UM                                3.9                      3.1
## 3 LM                                4                        4  
## 4 L                                 2.95                     4.1
## 5 <NA>                             NA                       NA  
## # i 40 more variables: median_expenditure_Europe <dbl>,
## #   `median_expenditure_Latin America and the Caribbean` <dbl>,
## #   `median_expenditure_Northern America` <dbl>,
## #   median_expenditure_Oceania <dbl>, median_expenditure_NA <dbl>,
## #   iqr_expenditure_Africa <dbl>, iqr_expenditure_Asia <dbl>,
## #   iqr_expenditure_Europe <dbl>,
## #   `iqr_expenditure_Latin America and the Caribbean` <dbl>, ...
\end{verbatim}

\begin{Shaded}
\begin{Highlighting}[]
\NormalTok{table\_test\_4 }\OtherTok{\textless{}{-}}\NormalTok{ table\_test }\SpecialCharTok{\%\textgreater{}\%}
  \FunctionTok{pivot\_longer}\NormalTok{(}
    \AttributeTok{cols =} \FunctionTok{starts\_with}\NormalTok{(}\StringTok{"median"}\NormalTok{),}
    \AttributeTok{names\_to =} \StringTok{"variable"}\NormalTok{,}
    \AttributeTok{values\_to =} \StringTok{"value"}
\NormalTok{  )}
\NormalTok{table\_test\_4}
\end{Highlighting}
\end{Shaded}

\begin{verbatim}
## # A tibble: 75 x 7
##    continent income_class iqr_expenditure iqr_unemployment iqr_migration
##    <chr>     <fct>                  <dbl>            <dbl>         <dbl>
##  1 Africa    H                       0                 0            0   
##  2 Africa    H                       0                 0            0   
##  3 Africa    H                       0                 0            0   
##  4 Africa    UM                      2.15             10.2          1.55
##  5 Africa    UM                      2.15             10.2          1.55
##  6 Africa    UM                      2.15             10.2          1.55
##  7 Africa    LM                      1.9              21.5          1.55
##  8 Africa    LM                      1.9              21.5          1.55
##  9 Africa    LM                      1.9              21.5          1.55
## 10 Africa    L                       2.55             14            3.05
## # i 65 more rows
## # i 2 more variables: variable <chr>, value <dbl>
\end{verbatim}

It is difficult to come up with anything meaningful beyond the basic
``table\_test'' structure. Any other way would deviate from tidy data
principles, be less efficient in terms of data representing the data,
and result in a significantly larger number of cells.

\begin{itemize}
\tightlist
\item
  Identify countries performing well in terms of both \textbf{youth
  unemployment} and \textbf{net migration rate} (top 25\% in net
  migration and bottom 25\% in youth unemployment within their
  continent).
\end{itemize}

\begin{Shaded}
\begin{Highlighting}[]
\NormalTok{top\_countries }\OtherTok{\textless{}{-}}\NormalTok{ df\_vars }\SpecialCharTok{\%\textgreater{}\%}
  \FunctionTok{group\_by}\NormalTok{(continent) }\SpecialCharTok{\%\textgreater{}\%}
  \FunctionTok{mutate}\NormalTok{(}
    \AttributeTok{migration\_q75 =} \FunctionTok{quantile}\NormalTok{(net\_migration\_rate, }\FloatTok{0.75}\NormalTok{, }\AttributeTok{na.rm =} \ConstantTok{TRUE}\NormalTok{),}
    \AttributeTok{unemployment\_q25 =} \FunctionTok{quantile}\NormalTok{(youth\_unempl\_rate, }\FloatTok{0.25}\NormalTok{, }\AttributeTok{na.rm =} \ConstantTok{TRUE}\NormalTok{)}
\NormalTok{  ) }\SpecialCharTok{\%\textgreater{}\%}
  \FunctionTok{filter}\NormalTok{(}
\NormalTok{    net\_migration\_rate }\SpecialCharTok{\textgreater{}=}\NormalTok{ migration\_q75,}
\NormalTok{    youth\_unempl\_rate }\SpecialCharTok{\textless{}=}\NormalTok{ unemployment\_q25}
\NormalTok{  ) }\SpecialCharTok{\%\textgreater{}\%}
  \FunctionTok{select}\NormalTok{(country, continent, net\_migration\_rate, youth\_unempl\_rate) }\SpecialCharTok{\%\textgreater{}\%}
  \FunctionTok{arrange}\NormalTok{(continent, net\_migration\_rate, }\FunctionTok{desc}\NormalTok{(youth\_unempl\_rate)) }\SpecialCharTok{\%\textgreater{}\%}
  \FunctionTok{ungroup}\NormalTok{()}

\NormalTok{top\_countries}
\end{Highlighting}
\end{Shaded}

\begin{verbatim}
## # A tibble: 17 x 4
##    country              continent           net_migration_rate youth_unempl_rate
##    <chr>                <chr>                            <dbl>             <dbl>
##  1 Togo                 Africa                             0                 3.9
##  2 Guinea               Africa                             0                 1  
##  3 Madagascar           Africa                             0                 1  
##  4 Benin                Africa                             0.3               5.6
##  5 Cote d'Ivoire        Africa                             1.2               5.5
##  6 Kazakhstan           Asia                               0.4               3.8
##  7 Macau                Asia                               3.3               5.3
##  8 Qatar                Asia                               6.5               0.4
##  9 United Arab Emirates Asia                               7.6               6.9
## 10 Bahrain              Asia                              10.6               5.3
## 11 Norway               Europe                             4                 9.7
## 12 Switzerland          Europe                             4.6               7.9
## 13 Malta                Europe                             6.6               9.1
## 14 Ecuador              Latin America and ~                0                 7.9
## 15 United States        Northern America                   3                 8.6
## 16 Papua New Guinea     Oceania                            0                 3.6
## 17 Palau                Oceania                            0.9               5.6
\end{verbatim}

\subsection{g. Conditional
Probabilities}\label{g.-conditional-probabilities}

Estimate the following based on the observed frequencies in the data:

\begin{itemize}
\item
  What is the (posterior or conditional) probability that a European
  country belongs to the high income group? What is the prior
  probability that a country belongs to the high income group?
\item
  Given a country has high youth unemployment (above \%25), what is the
  probability that it also has negative net migration?
\end{itemize}

\begin{Shaded}
\begin{Highlighting}[]
\NormalTok{p\_high\_given\_europe }\OtherTok{\textless{}{-}}\NormalTok{ df\_vars }\SpecialCharTok{\%\textgreater{}\%}
  \FunctionTok{filter}\NormalTok{(continent }\SpecialCharTok{==} \StringTok{"Europe"}\NormalTok{) }\SpecialCharTok{\%\textgreater{}\%}
  \FunctionTok{summarise}\NormalTok{(}\AttributeTok{prob =} \FunctionTok{mean}\NormalTok{(income\_class }\SpecialCharTok{==} \StringTok{"H"}\NormalTok{, }\AttributeTok{na.rm =} \ConstantTok{TRUE}\NormalTok{))}

\NormalTok{p\_high }\OtherTok{\textless{}{-}}\NormalTok{ df\_vars }\SpecialCharTok{\%\textgreater{}\%}
  \FunctionTok{summarise}\NormalTok{(}\AttributeTok{prob =} \FunctionTok{mean}\NormalTok{(income\_class }\SpecialCharTok{==} \StringTok{"H"}\NormalTok{, }\AttributeTok{na.rm =} \ConstantTok{TRUE}\NormalTok{))}

\NormalTok{p\_neg\_migration }\OtherTok{\textless{}{-}}\NormalTok{ df\_vars }\SpecialCharTok{\%\textgreater{}\%}
  \FunctionTok{filter}\NormalTok{(youth\_unempl\_rate }\SpecialCharTok{\textgreater{}} \DecValTok{25}\NormalTok{) }\SpecialCharTok{\%\textgreater{}\%}
  \FunctionTok{summarise}\NormalTok{(}\AttributeTok{prob =} \FunctionTok{mean}\NormalTok{(net\_migration\_rate }\SpecialCharTok{\textless{}} \DecValTok{0}\NormalTok{, }\AttributeTok{na.rm =} \ConstantTok{TRUE}\NormalTok{))}
\end{Highlighting}
\end{Shaded}

The chance that a European country is in the high-income group is about
76.09\%. The overall chance that any country belongs to the high-income
group, no matter the region, is 36.92\%. This large difference shows
that Europe has a higher number of high-income countries compared to the
global average. Being in Europe greatly raises the chances that a
country is classified as high-income. When having a high youth
unemployment rate the chances are with 61.22\% also high, that the
country has a negative migration rate.

\subsection{h. Simpson's Paradox
Analysis}\label{h.-simpsons-paradox-analysis}

Investigate whether an overall trend in youth unemployment rate in the
high and low income groups reverses when analyzed at the continent
level. E.g., does the youth unemployment rate appear lower in low-income
countries overall, but higher when controlling for continent? Explain
the results and possible reasons behind this paradox.

\begin{Shaded}
\begin{Highlighting}[]
\NormalTok{overall\_unemp }\OtherTok{\textless{}{-}}\NormalTok{ df\_vars }\SpecialCharTok{\%\textgreater{}\%}
  \FunctionTok{filter}\NormalTok{(income\_class }\SpecialCharTok{\%in\%} \FunctionTok{c}\NormalTok{(}\StringTok{"H"}\NormalTok{, }\StringTok{"L"}\NormalTok{)) }\SpecialCharTok{\%\textgreater{}\%}
  \FunctionTok{group\_by}\NormalTok{(income\_class) }\SpecialCharTok{\%\textgreater{}\%}
  \FunctionTok{summarise}\NormalTok{(}\AttributeTok{mean\_unemployment =} \FunctionTok{mean}\NormalTok{(youth\_unempl\_rate, }\AttributeTok{na.rm =} \ConstantTok{TRUE}\NormalTok{))}

\NormalTok{overall\_unemp}
\end{Highlighting}
\end{Shaded}

\begin{verbatim}
## # A tibble: 2 x 2
##   income_class mean_unemployment
##   <chr>                    <dbl>
## 1 H                         16.5
## 2 L                         14.7
\end{verbatim}

\begin{Shaded}
\begin{Highlighting}[]
\NormalTok{continent\_unemp }\OtherTok{\textless{}{-}}\NormalTok{ df\_vars }\SpecialCharTok{\%\textgreater{}\%}
  \FunctionTok{filter}\NormalTok{(income\_class }\SpecialCharTok{\%in\%} \FunctionTok{c}\NormalTok{(}\StringTok{"H"}\NormalTok{, }\StringTok{"L"}\NormalTok{)) }\SpecialCharTok{\%\textgreater{}\%}
  \FunctionTok{group\_by}\NormalTok{(continent, income\_class) }\SpecialCharTok{\%\textgreater{}\%}
  \FunctionTok{summarise}\NormalTok{(}\AttributeTok{mean\_unemployment =} \FunctionTok{mean}\NormalTok{(youth\_unempl\_rate, }\AttributeTok{na.rm =} \ConstantTok{TRUE}\NormalTok{)) }\SpecialCharTok{\%\textgreater{}\%}
  \FunctionTok{arrange}\NormalTok{(continent)}

\NormalTok{continent\_unemp}
\end{Highlighting}
\end{Shaded}

\begin{verbatim}
## # A tibble: 9 x 3
## # Groups:   continent [7]
##   continent                       income_class mean_unemployment
##   <chr>                           <chr>                    <dbl>
## 1 Africa                          H                         11.6
## 2 Africa                          L                         12.7
## 3 Asia                            H                         11.7
## 4 Asia                            L                         26.0
## 5 Europe                          H                         15.4
## 6 Latin America and the Caribbean H                         22.2
## 7 Northern America                H                         16.3
## 8 Oceania                         H                         25.7
## 9 <NA>                            H                        NaN
\end{verbatim}

Low-income countries have lower youth unemployment (14.68\%) than
high-income countries (16.54\%).

However, this pattern changes when looking at individual continents. In
Africa and Asia, low-income countries actually have higher unemployment
than high-income countries.

This is an example of Simpson's paradox, where a trend in the overall
data reverses when the data is analyzed in smaller groups. The reason is
that countries are unevenly distributed across continents, and regions
differ in their economic conditions.

\subsection{i. Data Export}\label{i.-data-export}

Export the final tidy dataset from e. as a \textbf{CSV} with:\texttt{;}
as a separator; \texttt{.} representing missing values; no row names
included. Upload the \texttt{.csv} to TUWEL, together with the
submission.

\begin{Shaded}
\begin{Highlighting}[]
\FunctionTok{write.table}\NormalTok{(df\_vars,}
            \AttributeTok{file =} \StringTok{"df\_vars.csv"}\NormalTok{,}
            \AttributeTok{sep =} \StringTok{";"}\NormalTok{,}
            \AttributeTok{na =} \StringTok{"."}\NormalTok{,}
            \AttributeTok{row.names =} \ConstantTok{FALSE}\NormalTok{)}
\end{Highlighting}
\end{Shaded}


\end{document}
